%\makeglossaries
\makenoidxglossaries

% Danh mục thuật ngữ và từ viết tắt
\newglossaryentry{AI}{
    type=\acronymtype,
    name={AI},
    description={Trí tuệ nhân tạo (Artificial Intelligence)},
    first={trí tuệ nhân tạo (Artificial Intelligence -- AI)}
}
\newglossaryentry{API}{
    type=\acronymtype,
    name={API},
    description={Giao diện lập trình ứng dụng (Application Programming Interface)},
    first={giao diện lập trình ứng dụng (Application Programming Interface -- API)}
}
\newglossaryentry{ASR}{
    type=\acronymtype,
    name={ASR},
    description={Nhận dạng tiếng nói tự động (Automatic Speech Recognition)},
    first={nhận dạng tiếng nói tự động (Automatic Speech Recognition -- ASR)}
}
\newglossaryentry{CLI}{
    type=\acronymtype,
    name={CLI},
    description={Giao diện dòng lệnh (Command-Line Interface)}
}
\newglossaryentry{CPU}{
    type=\acronymtype,
    name={CPU},
    description={Bộ xử lý trung tâm (Central Processing Unit)}
}
\newglossaryentry{CSDL}{
    type=\acronymtype,
    name={CSDL},
    description={Cơ sở dữ liệu}
}
\newglossaryentry{CSS}{
    type=\acronymtype,
    name={CSS},
    description={Ngôn ngữ định kiểu trang web (Cascading Style Sheets)}
}
\newglossaryentry{DI}{
    type=\acronymtype,
    name={DI},
    description={Tiêm phụ thuộc (Dependency Injection)}
}
\newglossaryentry{E2E}{
    type=\acronymtype,
    name={E2E},
    description={Kiểm thử đầu--cuối (End-to-End)}
}
\newglossaryentry{GPT}{
    type=\acronymtype,
    name={GPT},
    description={Mô hình ngôn ngữ lớn theo kiến trúc Generative Pre-trained Transformer của OpenAI}
}
\newglossaryentry{GPU}{
    type=\acronymtype,
    name={GPU},
    description={Bộ xử lý đồ họa (Graphics Processing Unit)}
}
\newglossaryentry{HTTP}{
    type=\acronymtype,
    name={HTTP},
    description={Giao thức truyền siêu văn bản (HyperText Transfer Protocol)}
}
\newglossaryentry{HTTPS}{
    type=\acronymtype,
    name={HTTPS},
    description={Giao thức truyền siêu văn bản bảo mật (HyperText Transfer Protocol Secure)}
}
\newglossaryentry{IRP}{
    type=\acronymtype,
    name={IRP},
    description={Thực hành truy hồi xen kẽ (Interpolated Retrieval Practice)}
}
\newglossaryentry{JSON}{
    type=\acronymtype,
    name={JSON},
    description={Định dạng trao đổi dữ liệu JavaScript Object Notation}
}
\newglossaryentry{JWT}{
    type=\acronymtype,
    name={JWT},
    description={Chuỗi xác thực JSON Web Token},
    first={JSON Web Token (JWT)}
}
\newglossaryentry{LLM}{
    type=\acronymtype,
    name={LLM},
    description={Mô hình ngôn ngữ lớn (Large Language Model)},
    first={mô hình ngôn ngữ lớn (Large Language Model -- LLM)}
}
\newglossaryentry{MCQ}{
    type=\acronymtype,
    name={MCQ},
    description={Câu hỏi trắc nghiệm nhiều lựa chọn (Multiple Choice Question)},
    first={câu hỏi trắc nghiệm nhiều lựa chọn (Multiple Choice Question -- MCQ)}
}
\newglossaryentry{MOOC}{
    type=\acronymtype,
    name={MOOC},
    description={Khóa học trực tuyến mở đại trà (Massive Open Online Course)}
}
\newglossaryentry{MP4}{
    type=\acronymtype,
    name={MP4},
    description={Định dạng tệp video MPEG-4}
}
\newglossaryentry{RAM}{
    type=\acronymtype,
    name={RAM},
    description={Bộ nhớ truy cập ngẫu nhiên (Random Access Memory)}
}
\newglossaryentry{REST}{
    type=\acronymtype,
    name={REST},
    description={Kiểu kiến trúc dịch vụ web Representational State Transfer}
}
\newglossaryentry{RPC}{
    type=\acronymtype,
    name={RPC},
    description={Gọi thủ tục từ xa (Remote Procedure Call)},
    first={gọi thủ tục từ xa (Remote Procedure Call -- RPC)}
}
\newglossaryentry{RTF}{
    type=\acronymtype,
    name={RTF},
    description={Hệ số thời gian thực (Real-Time Factor)}
}
\newglossaryentry{S3}{
    type=\acronymtype,
    name={S3},
    description={Giao thức lưu trữ đối tượng tương thích Amazon Simple Storage Service},
    first={giao thức lưu trữ đối tượng tương thích Amazon Simple Storage Service (S3)}
}
\newglossaryentry{SQL}{
    type=\acronymtype,
    name={SQL},
    description={Ngôn ngữ truy vấn có cấu trúc (Structured Query Language)}
}
\newglossaryentry{SSR}{
    type=\acronymtype,
    name={SSR},
    description={Kết xuất phía máy chủ (Server-Side Rendering)}
}
\newglossaryentry{STT}{
    type=\acronymtype,
    name={STT},
    description={Chuyển tiếng nói thành văn bản (Speech-to-Text)},
    first={chuyển tiếng nói thành văn bản (Speech-to-Text -- STT)}
}
\newglossaryentry{TLS}{
    type=\acronymtype,
    name={TLS},
    description={Giao thức bảo mật tầng truyền tải (Transport Layer Security)}
}
\newglossaryentry{TOML}{
    type=\acronymtype,
    name={TOML},
    description={Định dạng tệp cấu hình Tom's Obvious, Minimal Language}
}
\newglossaryentry{TTS}{
    type=\acronymtype,
    name={TTS},
    description={Tổng hợp tiếng nói từ văn bản (Text-to-Speech)},
    first={tổng hợp tiếng nói từ văn bản (Text-to-Speech -- TTS)}
}
\newglossaryentry{UI}{
    type=\acronymtype,
    name={UI},
    description={Giao diện người dùng (User Interface)}
}
\newglossaryentry{URL}{
    type=\acronymtype,
    name={URL},
    description={Định vị tài nguyên thống nhất (Uniform Resource Locator)}
}
\newglossaryentry{UUID}{
    type=\acronymtype,
    name={UUID},
    description={Định danh duy nhất toàn cục (Universally Unique Identifier)}
}
\newglossaryentry{VRAM}{
    type=\acronymtype,
    name={VRAM},
    description={Bộ nhớ truy cập ngẫu nhiên của card đồ họa (Video RAM)}
}
\newglossaryentry{WAV}{
    type=\acronymtype,
    name={WAV},
    description={Định dạng tệp âm thanh dạng sóng (Waveform Audio File Format)}
}
\newglossaryentry{WER}{
    type=\acronymtype,
    name={WER},
    description={Tỷ lệ lỗi từ (Word Error Rate)}
}
